% !TEX program = xelatex
% 编译方式：在本目录执行两遍 xelatex（详见 docs/_manual_common/README.md）。
% 源码依据：单头文件 include/hacdcpf/analysis/harmonics_power_flow.hpp 与
%           src/harmonics_power_flow/harmonics_power_flow.cpp 源码，公式与参数以代码为准。
\documentclass[UTF8,fontset=fandol,a4paper,11pt]{ctexart}
\usepackage{hysim_manual}

\title{\textbf{交直流混合高弹性能源电力系统仿真平台}\\[0.5em]
       {\Large 谐波潮流技术手册}\\[0.3em]
       {\large 数学模型、接口参数、验证校核与工业级评价}}
\author{HySim-XJTU-HRPES（西安交通大学 HRPES 团队）}
\date{\today}

\begin{document}
\maketitle
\begin{abstract}
\noindent 本手册依据 HySim-XJTU-HRPES 平台谐波潮流模块整理。该模块以单个头文件
\srcpath{include/hacdcpf/analysis/harmonics\_power\_flow.hpp} 与实现
\srcpath{src/harmonics\_power\_flow/harmonics\_power\_flow.cpp} 提供（注意：并无
\srcpath{include/hacdcpf/harmonics\_power\_flow/} 目录），命名空间 \texttt{hacdcpf::harmonics}。
手册覆盖频域“谐波穿透”线性直接解、牛顿谐波平衡变体、三相 abc 域求解与 NIC 交直流桥接，以及
IEEE 519-2014 / GB/T 14549-1993 合规校核。顶层暂无权威 Markdown 文档，头文件与测试即权威来源；
直流纹波网络按设计为纯阻性，基频未收敛时回退到存储/额定电压并在 \fld{model\_limitations} 中标注。
版式与《碳流追踪与碳排放核算技术手册》一致。
\end{abstract}

\tableofcontents
\newpage

\section{阅读指南与约定}
\subsection{数据来源}
\begin{itemize}
  \item \srcpath{include/hacdcpf/analysis/harmonics\_power\_flow.hpp} —— 全部接口、选项与结果结构
        （单头文件）；
  \item \srcpath{src/harmonics\_power\_flow/harmonics\_power\_flow.cpp} —— 求解实现；
  \item \srcpath{tests/test\_harmonics\_power\_flow.cpp} 与 \srcpath{tools/harmonics\_validation/}
        —— 验证与跨引擎校核依据；
  \item 归档理论 \srcpath{docs/archive/theory/harmonics/}（\srcpath{harmonic\_general.md}、
        \srcpath{harmonic\_nic.md}、\srcpath{harmonic\_three\_phase.md} 等）。
\end{itemize}
命名空间为 \texttt{hacdcpf::harmonics}。顶层无权威 Markdown，契约以头文件与测试为准。

\subsection{单位制与索引空间}
谐波次 $h$ 为整数（交流缺省奇次、直流缺省偶次），电压以标幺或母线额定折算，THD 以百分数表示。
交直流结果分列：\fld{ac\_bus\_results}/\fld{dc\_bus\_results}（\texttt{HarmonicBusResult}）、
\fld{ac\_branch\_flows}；每母线以 \fld{v\_by\_order}（次$\to$复数）、\fld{thd\_pct}、\fld{v\_fund\_pu}
表述。索引空间与基频潮流的交/直流母线一致。

\subsection{诚实口径}
本模块为“Level 1”线性模型：每个谐波次一次复数稀疏分解。直流纹波网络按设计为纯阻性
（DCBranch 无电感），除非 \texttt{DCRippleModel} 覆盖引入 \fld{x\_pu}/\fld{b\_pu}。负荷阻性部分
$P/V^2$ 与频率无关，不随集肤效应缩放（仅线路/变压器/电源串联 R 频变）。三相
\srcpath{solve\_harmonic\_power\_flow\_3ph} 求解不含直流子系统（仅交流），此时
\fld{auto\_nic\_from\_vscs} 被忽略。请求但未收敛的基频潮流回退到存储/额定电压
（\fld{used\_stored\_operating\_point}=true，并列入 \fld{model\_limitations}）。

\section{功能定位与架构}
\subsection{公共入口族}
\begin{center}\footnotesize
\begin{tabular}{@{}L{6.6cm}L{3.0cm}L{5.4cm}@{}}
\toprule
入口 & 定义位置 & 说明\\
\midrule
\srcpath{solve\_harmonic\_power\_flow} & analysis/harmonics\_power\_flow.hpp & 单相/正序交直流线性直接节点解（Level 1）\\
\srcpath{solve\_harmonic\_power\_flow\_newton} 等 & 同上 & 牛顿谐波平衡变体（含 \_real/\_coupled）\\
\srcpath{solve\_harmonic\_power\_flow\_3ph} & 同上 & 三相 abc 域求解（仅交流）\\
\srcpath{solve\_harmonic\_power\_flow\_3ph\_hybrid} & 同上 & 三相 + NIC 交直流桥接\\
\srcpath{solve\_harmonic\_power\_flow\_hybrid\_newton} & 同上 & 交直流耦合牛顿求解\\
\srcpath{check\_harmonic\_limits} / \srcpath{harmonic\_voltage\_limits} & 同上 & IEEE 519 / GB/T 14549 合规校核\\
\srcpath{default\_six\_pulse\_ac\_spectrum} / \srcpath{default\_dc\_ripple\_spectrum} & 同上 & 6 脉动交流/直流纹波缺省谱\\
\bottomrule
\end{tabular}
\end{center}

\subsection{关键数据结构}
\compmeta{HPFOptions}{include/hacdcpf/analysis/harmonics\_power\_flow.hpp}
主要字段：\fld{ac\_orders}（缺省 5, 7, 11, 13, 17, 19, 23, 25）、\fld{dc\_orders}
（缺省 2, 6, 12, 18, 24）、\fld{run\_base\_power\_flow}、\fld{include\_load\_impedance}、
\fld{default\_source\_xpp\_pu}（$0.2$）、\fld{skin\_effect}（\texttt{SkinEffectModel}）、
\fld{auto\_nic\_from\_vscs}、\fld{newton\_max\_iter}/\fld{newton\_tol}。

\compmeta{HPFResult}{include/hacdcpf/analysis/harmonics\_power\_flow.hpp}
主要字段：\fld{ok}、\fld{ac\_bus\_results}/\fld{dc\_bus\_results}、\fld{ac\_branch\_flows}、
\fld{base\_pf\_converged}、\fld{used\_stored\_operating\_point}、\fld{model\_limitations}、
\fld{max\_ac\_thd\_pct}。两端口 \texttt{HarmonicNIC} 含 \fld{ac\_port}/\fld{dc\_port}
（\texttt{PortBehavior}）、\fld{ac\_spectrum}/\fld{dc\_spectrum}、\fld{y\_out\_ac}；单端口
\texttt{HarmonicCurrentSource} 给出诺顿谱。合规结构 \texttt{HarmonicComplianceReport}/
\texttt{HarmonicLimitCheck} 含 \fld{thd\_ok}、\fld{worst\_ihd\_order}、\fld{compliant} 与
\texttt{HarmonicStandard}。

\section{核心数学模型}
频域“谐波穿透”把各谐波次视为独立的复数节点问题，资源以诺顿谐波电流源表示。

\subsection{频域节点方程}
对每个谐波次 $h$ 组装复数节点导纳并直接求解
\begin{equation}
\mathbf{Y}(h)\,\mathbf{V}(h) = \mathbf{I}(h),\qquad h\in\{5,7,11,\dots\},
\end{equation}
每次仅一次复数稀疏分解（Level 1，逐次线性）。

\subsection{NIC 一致注入}
两端口 NIC 由同一基频通过功率给出自洽的交/直流谱
\begin{equation}
I_{ac,1} = \frac{\overline{S_{ac}}}{\overline{V_{ac,1}}},\qquad
I_{dc,0} = \frac{P_{dc}}{V_{dc,0}}.
\end{equation}

\subsection{支路与负荷谐波阻抗}
交流支路按
\begin{equation}
Z_{ij}(h) = r_{ij} + \mathrm{j}\,h\,x_{ij},\qquad
y^{\text{sh}}_{ij}(h) = \mathrm{j}\,h\,b_{ij}
\end{equation}
缩放，负荷以 CIGRE/IEC 并联 $R\parallel\mathrm{j}X$ 建模。

\subsection{三相序分量}
三相路径用
\begin{equation}
\mathbf{Z}_{abc}(h) = \mathbf{A}\,\mathrm{diag}(z_0,z_1,z_1)\,\mathbf{A}^{-1},
\end{equation}
按 $h \bmod 3$ 自动指派零/正/负序。

\subsection{THD 与合规}
电压总谐波畸变率
\begin{equation}
\mathrm{THD}_v = \frac{\sqrt{\sum_{h>1}\lvert V_h\rvert^2}}{\lvert V_1\rvert}\times 100\%,
\end{equation}
并按 IEEE 519-2014 / GB/T 14549-1993 判定 \fld{thd\_ok} 与 \fld{compliant}。

\section{接口与参数}
\begin{paramtable}
\fld{ac\_orders} & — & — & 奇次为主 & 见说明 & 交流谐波次集，缺省 5, 7, 11, 13, 17, 19, 23, 25\\
\fld{dc\_orders} & — & — & 偶次为主 & 见说明 & 直流纹波次集，缺省 2, 6, 12, 18, 24\\
\fld{run\_base\_power\_flow} & — & — & 布尔 & — & 先解基频潮流以取工作点\\
\fld{include\_load\_impedance} & — & — & 布尔 & — & 负荷并入谐波阻抗（并联 R//jX）\\
\fld{default\_source\_xpp\_pu} & $x''$ & pu & — & $0.2$ & 电源次暂态电抗缺省值\\
\fld{skin\_effect} & — & — & 枚举 & — & 集肤效应模型（仅作用于串联 R）\\
\fld{auto\_nic\_from\_vscs} & — & — & 布尔 & — & 由 VSC 自动生成两端口 NIC\\
\fld{newton\_max\_iter} & — & — & 整数 & — & 牛顿谐波平衡最大迭代次数\\
\fld{newton\_tol} & — & — & 小正数 & — & 牛顿谐波平衡收敛容差\\
\fld{max\_ac\_thd\_pct} & THD & \% & — & — & 结果：全网最大交流电压 THD\\
\fld{used\_stored\_operating\_point} & — & — & 布尔 & — & 结果：基频未收敛时回退标志\\
\end{paramtable}

\section{验证与边界局限}
\subsection{验证与校核}
注册测试 \srcpath{test\_harmonics\_power\_flow}（\srcpath{tests/test\_harmonics\_power\_flow.cpp}）；
ctest 项 \srcpath{harmonics\_cross\_engine\_matrix}（标签 \texttt{harmonics;validation;opendss}）；
校核驱动 \srcpath{validate\_harmonics\_xref}、\srcpath{validate\_harmonics\_ieee13} 与目标
\srcpath{validate\_harmonics\_device\_opendss}（\srcpath{tools/harmonics\_validation/}）。

\subsection{边界与局限（如实标注）}
\begin{itemize}
  \item “Level 1”线性模型；直流纹波网络按设计纯阻性（DCBranch 无电感），除非
        \texttt{DCRippleModel} 覆盖引入 \fld{x\_pu}/\fld{b\_pu}；
  \item 负荷阻性部分 $P/V^2$ 频率无关，不随集肤效应缩放（仅串联 R 频变）；
  \item 三相 \srcpath{solve\_harmonic\_power\_flow\_3ph} 不建模直流子系统（仅交流），
        \fld{auto\_nic\_from\_vscs} 在其中被忽略；
  \item 基频请求但未收敛时回退存储/额定电压（\fld{used\_stored\_operating\_point}=true，列入
        \fld{model\_limitations}）；
  \item 顶层无权威 Markdown，头文件与测试即权威；归档理论见
        \srcpath{docs/archive/theory/harmonics/}。
\end{itemize}
\section{跨模块工业级技术评价基线}

本章规定本手册所述模块进入工程算例、批量研究或生产服务前必须共同满足的评价基线。
具体模型公式、变量、约束和专用算法以正文相应章节为准；本章不扩大源码已经实现的范围。

\subsection{输入、输出、边界条件与数据结构审计}
输入审计应逐项检查有限性、单位、上下界、枚举值和引用完整性。系统对象必须先按所需层级完成
校验；AC 与 DC 母线采用域限定映射，组件稳定 \texttt{index}、向量位置和临时图索引不得混用。
输出审计必须记录单位、索引空间、模型范围、收敛或可行状态、后端、容差、迭代/节点计数和告警。
空系统、空时域、孤岛、重复 ID、零容量、退化约束与不可用可选后端均属于边界测试集。

\subsection{工程假设与数值稳定性分析}
模型使用的平衡、线性化、稳态、独立性、概率分布、时间离散或网络等值假设必须与应用场景匹配。
数值评价至少包含变量缩放、绝对/相对残差、可行性违反、条件数或枢轴质量、步长/阻尼、迭代停滞
和随机种子复现性。若采用正则化、平滑、回退、松弛或近似，应同时报告触发条件、参数值、对主指标
的敏感性以及是否改变模型口径；不得用“已返回结果”替代收敛或有效性证明。

\subsection{复杂度与资源分析}
令 $n_x$、$n_c$、$n_t$、$n_s$、$|V|$、$|E|$ 分别表示变量、约束、时段、场景、节点和边数。
报告应分别给出建模、求解、后处理的时间与内存。图遍历通常为 $O(|V|+|E|)$；逐时段/逐场景
流程至少线性增长为 $O(n_t n_s)$ 次子问题调用；稠密线性代数最坏可达 $O(n_x^3)$，稀疏方法则应
报告结构非零数、填充率和因子分解峰值内存；MILP/NLP 不给出虚假的多项式最坏界，而报告节点数、
间隙、时限和实测扩展曲线。复杂度结论必须基于固定硬件、固定后端和可复现算例族。

\subsection{异常工况处理}
非法输入应在进入求解器前返回结构化错误；不可行、无界、不收敛、时限、内存不足和后端缺失必须
相互区分。部分结果只有在对应 \fld{ValidityFlags}、\fld{model\_scope}、
\fld{model\_limitations} 或等价字段允许时才可使用。异常路径不得静默丢弃 AC/DC 资产、制造平衡
节点、把 NaN 改写为零或把近似解标为全局最优；系统原始对象应保持可恢复和可重试。

\subsection{与其他模块的接口关系}
接口链路遵循“工程语义模型 $\to$ 校验 $\to$ 投影/装配 $\to$ 求解或分析 $\to$ 结果回投影”。
跨模块传递时保留稳定 ID、单位、符号、时间基准和场景身份；消费潮流、OPF、动态或可靠性结果前，
先检查其收敛、覆盖范围和有效性标志。HTTP/JSON/Python 层只做语义保持适配，不重新解释求解结果。

\subsection{测试与验证建议}
最低验证矩阵包括：手算小例、标准算例、边界/异常输入、随机性质测试、算法或后端交叉校核、
序列化往返、稳定 ID 回投影、AC/DC 同号母线、重复运行确定性和规模扩展测试。数值算法还应加入
残差独立重算、雅可比有限差分或自动微分校核；近似模型应与高保真模型比较并给出误差分布，而非
只报告单个平均值。回归报告必须列明构建配置、算例版本、容差、通过/失败/跳过数及未验证范围。

\subsection{工业级评估指标}
\begin{itemize}
  \item \textbf{正确性}：守恒残差、约束最大违反、解析/基准误差、结果回投影一致性；
  \item \textbf{鲁棒性}：收敛率、不可行识别率、异常分类准确率、扰动敏感性与随机复现率；
  \item \textbf{性能}：端到端时延、吞吐量、峰值内存、并行效率与规模增长斜率；
  \item \textbf{可追溯性}：输入模式版本、稳定 ID 覆盖率、模型范围/限制字段完整率、告警可定位率；
  \item \textbf{工程有效性}：与标准、外部引擎或现场数据的偏差，以及误判/漏判对业务指标的影响。
\end{itemize}
指标应预先给出验收阈值和失败处置；未达到阈值时只能标记为实验性或受限适用。

\subsection{适用范围、局限性与后续改进}
本手册只评价当前源码覆盖的模型、后端和数据结构，不对未建模物理过程、未安装求解器或未执行的
外部验证作保证。后续改进应按“缺口 $\to$ 可量化指标 $\to$ 源码/测试变更 $\to$ 独立验证”闭环，
优先消除会造成错误结论的语义缺口，其次提升数值鲁棒性与性能，最后扩展模型范围。


\end{document}
